\chapter{Geometry, Color, Curves, and C\texttt{++} Value Layout}
\label{ch:geometry-color-layout}

Chapter~\ref{ch:math-conventions} gives the method-level reference for CNA's math types.  This
chapter takes the complementary view: which geometric conventions cross subsystem boundaries,
which C\texttt{++} object layouts differ from XNA value types, and which apparently harmless casts
or degeneracies become observable defects.

\section{The value-type resemblance has a hard limit}

XNA presents vectors, colours, planes, bounds, and curve keys as compact managed value types.
CNA preserves their names and most operations, but it cannot preserve that representation merely
by spelling the same fields in C\texttt{++}.  Several public types inherit polymorphic interfaces:

\begin{center}
\small
\begin{tabularx}{\textwidth}{>{\raggedright\arraybackslash}p{3.1cm}>{\raggedright\arraybackslash}p{4.1cm}X}
\toprule
Type family & Why it is polymorphic & Practical consequence \\
\midrule
\cnaclass{Color} & Implements \texttt{IPackedVectorT<}\newline\texttt{UInt32>}%
  \index{IPackedVectorT<UInt32>}, whose base has virtual
  operations and a virtual destructor. & The object has a vptr; it is not four raw RGBA bytes. \\
Built-in vertex types & Implement \cnaclass{IVertexType}. & A public vertex object cannot be
  uploaded by copying \texttt{sizeof(T)} bytes. \\
\cnaclass{BoundingBox}, \cnaclass{BoundingSphere},\newline
\cnaclass{BoundingFrustum}, \cnaclass{CurveKey} & Implement \cnaclass{IEquatable<T>}. &
  Their field list is not their complete
  object representation. \\
Plain math types & \cnaclass{Vector2}, \cnaclass{Vector3}, \cnaclass{Vector4}, \cnaclass{Matrix},
  \cnaclass{Quaternion}, \cnaclass{Plane}, \cnaclass{Ray}, \cnaclass{Point}, and
  \cnaclass{Rectangle} have no base class. & Their layout is simpler, but padding and C\texttt{++}
  object rules still belong to the ABI, not to the XNA contract. \\
\bottomrule
\end{tabularx}
\end{center}

At the pinned revision, \texttt{sizeof(Color) == 24}.  The single packed value is still
\texttt{AABBGGRR} on a little-endian machine, yet the object begins with a vptr.  Reinterpreting a
\texttt{Color*} as pixel bytes would overwrite or read the virtual-table pointer.  CNA's own
\texttt{GetBackBufferData<Color>} path therefore reads native bytes into a temporary byte vector
and constructs each result as \texttt{Color(r,g,b,a)}.  \cnaclass{Texture3D} and
\cnaclass{TextureCube} use the same conversion rule.

This is not theoretical portability advice.  A regression test explicitly rejects the old raw
cast, and the graphics module documents the exact 24-byte size where it converts texture data.
Application code should use component access, \texttt{PackedValue}, or the typed texture APIs; it
should never serialize a public math object with \texttt{memcpy(sizeof(T))} unless the file format
is deliberately CNA-ABI-specific.

\section{GPU streams are a second, plain-data representation}

The graphics core solves the polymorphic-vertex problem with a parallel internal stream layer in
\path{modules/graphics/include/CNA/Internal/Graphics/BuiltInVertexStreams.hpp}.  These structs have
no virtual base, and both their total sizes and member offsets are compile-time assertions:

\begin{center}
\small
\begin{tabular}{lr}
\toprule
Stream representation & Bytes \\
\midrule
\texttt{PositionColorStream} & 16 \\
\texttt{PositionTextureStream} & 20 \\
\texttt{PositionColorTextureStream} & 24 \\
\texttt{PositionNormalTextureStream} & 32 \\
\texttt{PositionNormalTangentTextureStream} & 48 \\
\texttt{PositionNormalTextureSkinnedStream} & 52 \\
\texttt{PositionNormalTangentTextureSkinnedStream} & 68 \\
\bottomrule
\end{tabular}
\end{center}

Every built-in \cnaclass{VertexDeclaration} takes its stride from the matching stream struct, not
from the public vertex class.  Conversion is therefore an explicit semantic operation: positions
and texture coordinates are copied as floats, while \cnaclass{Color} contributes only its packed
four-byte payload.  Chapter~\ref{ch:vertex-streams-capabilities} follows these bytes through binding,
multi-stream normalization, and renderer capability gates.

\section{Planes and half-spaces}

CNA uses the plane equation

\[
  \mathbf{n}\cdot\mathbf{x} + D = 0.
\]

The three-point constructor is the only constructor that normalizes automatically and computes
\(D=-\mathbf{n}\cdot\mathbf{a}\).  \texttt{DotCoordinate} returns the signed expression above.
Operationally, \cnaclass{PlaneIntersectionType::Front} is the positive-normal side and
\texttt{Back} is the negative side; the enum's own prose comments reverse those meanings, while
\cnaclass{Plane}, \cnaclass{BoundingBox}, and \cnaclass{BoundingSphere} consistently implement the
positive-side interpretation.

\subsection{The inverse-transpose path is presently wrong}

Transforming a plane should multiply it by the inverse transpose of the point transform.  CNA's
\texttt{Plane::}\allowbreak\texttt{Transform}\index{Plane::Transform} intends to do that, but calls
the out-parameter transpose as
\texttt{Matrix::}\allowbreak\texttt{Transpose(m, m)}.  That overload writes fields while it is
still reading the same
object, so it is not alias-safe.  For a non-symmetric inverse it progressively mixes old and new
fields and produces a symmetrization rather than a transpose.

A minimal counterexample is the plane \(y=1\) transformed by a pure translation along X.  The
correct plane is unchanged.  The aliased implementation instead produces a normal proportional to
\((1,1,0)\), silently rotating it.  Both existing transform tests use the identity matrix, whose
inverse is symmetric, so they cannot expose the error.  This is recorded as CNA-BUG-001; portable
application code should avoid \cnaclass{Plane::Transform} for nontrivial matrices until the source
uses a separate transpose destination.

\section{Bounds encode several non-obvious conventions}

\cnaclass{BoundingFrustum} independently confirms CNA's Direct3D clip volume.  Its extracted planes
describe \(-w\le x,y\le w\) and \(0\le z\le w\), and their normals point outward.  Consequently a
point on the positive side of any plane is outside.  This agrees with the projection builders in
Chapter~\ref{ch:math-conventions}; it is a second derivation, not a repeated assertion.

The three bounding types do not share one quality level:

\begin{itemize}
  \item \cnaclass{BoundingBox::GetCorners} preserves XNA's fixed eight-corner order.  Its
        containment tests use half-space and min/max logic, but
        \texttt{Contains(BoundingFrustum)} retains legacy reversed-question behavior.
  \item \cnaclass{BoundingSphere} builds point sets with \texttt{CreateFromPoints}, using a
        Ritter-style expansion.\index{BoundingSphere::CreateFromPoints} Under a non-uniform
        transform, the radius is multiplied by the largest row-vector scale.  Its frustum
        \texttt{Contains} path never returns \texttt{Disjoint} because the
        required minimum-distance accumulator is never populated.
  \item \cnaclass{BoundingFrustum} caches six planes and eight corners.  Its
        \texttt{Intersects(Ray)} main case throws \cnaclass{NotImplementedException}; source and
        test absence must not be mistaken for a completed geometric query.
\end{itemize}

Those gaps contain no \texttt{TODO} marker.  They are a useful reminder that marker searches find
intent labels, not semantic incompleteness.

\section{Rectangle, Point, and edge ownership}

\cnaclass{Rectangle::Contains} is half-open: the left and top edges belong to the rectangle, while
the right and bottom edges do not.  \texttt{Intersects} requires positive overlap, so rectangles
that merely touch at an edge are not intersecting.  \texttt{Center} uses integer arithmetic, and
\texttt{IsEmpty} is true only when all four fields are zero; a zero-width rectangle at a nonzero
position is not the singleton \texttt{Empty} value.

\cnaclass{Point} is simply two 32-bit integers.  It deliberately has no implicit
\cnaclass{Vector2} bridge.  Make rounding policy visible when converting a continuous coordinate to
a pixel or cell: truncation, floor, and nearest are different behaviors for negative values.

\section{Color is packed, named, and polymorphic}

The public surface contains 141 named colours: 140 opaque values plus \texttt{Transparent}.  There
are 139 distinct packed values because \texttt{Aqua}/\texttt{Cyan} and
\texttt{Fuchsia}/\texttt{Magenta} are aliases.  Three counts can therefore all appear plausible;
only a count that names declarations, opacity, or distinct packed values is meaningful.

Construction from float vectors clamps components to \([0,1]\), scales to \([0,255]\), and packs
the four bytes.  Integer constructors clamp to \([0,255]\).  Multiplication scales alpha as well as
RGB.  These behavioral facts are independent of the 24-byte C\texttt{++} object layout: packing
defines the value, not permission to treat the entire object as packed storage.

\section{Curves: ordered keys and explicit extrapolation}

\cnaclass{Curve} evaluates cubic Hermite segments between sorted \cnaclass{CurveKey} entries.
Duplicate positions are allowed; insertion and \texttt{setItemProperty} preserve sorted order and
can reposition a key.  \cnaclass{CurveContinuity::Step} selects the left value until the segment
endpoint, while smooth keys use their incoming and outgoing tangents.

Five loop modes govern positions outside the key range.  \texttt{Constant} and \texttt{Linear}
hold or extrapolate an endpoint; \texttt{Cycle}, \texttt{CycleOffset}, and \texttt{Oscillate}
repeat a span.  The repeating modes depend on a stable cycle count and on a nonzero range.  CNA
guards some degenerate spans, but its tangent computation
uses different near-zero tests for incoming and outgoing sides.  Treat a duplicate-position curve
as an authored edge case and test it with the same evaluator that will ship.

Curves also cross the content boundary.  CNJ can carry a self-contained Curve payload, and the XNB
reader has a concrete \texttt{CurveReader}; a successful parse proves field recovery, while the
runtime evaluator and loop behavior remain separate claims.

\section{Exceptions are part of the port}

Sibling operations do not consistently choose the same exception family.  With empty input,
\texttt{CreateFromPoints}\index{BoundingBox::CreateFromPoints}%
\index{BoundingSphere::CreateFromPoints} throws different exceptions.  \cnaclass{BoundingBox}
uses \cnaclass{ArgumentException}; \cnaclass{BoundingSphere} uses the standard invalid-argument
exception.

An undersized corner destination
throws \cnaclass{ArgumentOutOfRangeException} for the box and \texttt{std::out\_of\_range} for the
frustum.  Existing tests encode these differences, so a catch block written for one type is not a
generic geometry policy.

Degenerate numeric inputs are usually quieter: normalizing a zero vector or quaternion and
inverting a singular matrix can yield NaN or infinity.  The module's strong argument checks are
concentrated in projection factories.  Validate user-authored axes, scales, planes, and key spans at
the application boundary rather than expecting every math primitive to reject them.

\section{What the math tests prove}

The pinned math module contains 818 statically counted GoogleTest definitions.  They use
\texttt{EXPECT\_FLOAT\_EQ} or explicit tolerances of \(10^{-4}\), \(10^{-5}\), and \(10^{-6}\);
none uses captured XNA or MonoGame numeric oracle values.  The suite is strong on API wiring and
many analytic identities, but it does not independently pin matrix/quaternion composition order,
the \([0,1]\) depth mapping, corner order, singular inputs, or the non-identity plane transform.

For a port, pair these types with tests derived from the original game's observable behavior:
screen-space points, collision classifications, curve samples, or captured reference values.  A
familiar type name establishes lineage.  It does not erase C\texttt{++} layout, exception, and
evidence boundaries.
